\documentclass[11pt,a4paper]{article}
\usepackage[margin=2.3cm]{geometry}
\usepackage{amsmath,amssymb,amsthm,mathtools,bm}
\usepackage{booktabs}
\usepackage[hidelinks]{hyperref}
\usepackage{enumitem}
\setlist{itemsep=2pt,topsep=4pt}

\newtheorem*{remark}{Remark}
\newtheorem*{prop}{Proposition}
\newcommand{\UB}{\mathrm{UB}}
\newcommand{\LB}{\mathrm{LB}}
\newcommand{\VWAP}{\mathrm{VWAP}}
\newcommand{\E}{\mathbb{E}}
\newcommand{\1}{\mathbf{1}}

\title{Intraday Momentum on SPY\\[2pt]\large A Mathematical Description of the Strategy and its Evaluation}
\author{Emil Ohligs}
\date{Based on Zarattini, Aziz \& Barbon (2024), \emph{Beat the Market}}

\begin{document}
\maketitle

\noindent This note states every step of the strategy as a formula, in the same order as the code in
\texttt{src/intraday\_momentum/}. Each section names the module that implements it.

\tableofcontents

%=======================================================================
\section{Notation and information structure}
%=======================================================================

\paragraph{Days and minutes.}
Trading days are indexed by $t = 1,\dots,T$. The regular session (09:30--16:00 ET) is divided into
$N = 390$ one-minute bars, indexed by $j = 0,\dots,N-1$. Bar $j$ covers the interval
$[\,09{:}30 + j,\; 09{:}30 + j + 1)$. On half days the session ends after $N_t < N$ bars;
in general $N_t \le N$ is the number of bars on day $t$ (\texttt{data.py}).

\paragraph{Prices.} For each day and bar we observe open, high, low, close and volume
$(O_{t,j}, H_{t,j}, L_{t,j}, C_{t,j}, V_{t,j})$. We write
\[
  P_{t,k} := C_{t,k-1}, \qquad k = 1,\dots,N_t,
\]
for the price observed at minute $k$ after the open (the close of the bar that ends at $09{:}30+k$).
The day's open and close are
\[
  O_t := O_{t,0}, \qquad C_t := P_{t,N_t}.
\]

\paragraph{Filtration.} Let $\mathcal F_{t,k}$ be the information available at minute $k$ of day $t$:
all bars of days $1,\dots,t-1$ and bars $0,\dots,k-1$ of day $t$.
A backtest is free of look-ahead bias if and only if every decision taken at $(t,k)$ is
$\mathcal F_{t,k}$-measurable and executed at a price that is realised \emph{after} $(t,k)$.
Every quantity below is marked with the information set it depends on.

%=======================================================================
\section{The Noise Area (\texttt{signals.py})}
%=======================================================================

\subsection{Average move from the open}
On every past day $t-i$ we measure the absolute relative move from the open after $k$ minutes:
\begin{equation}
  m_{t-i,k} = \left| \frac{P_{t-i,k}}{O_{t-i}} - 1 \right| .
\end{equation}
The \emph{expected size of a normal move} at minute $k$ is estimated by the mean over the
previous $n = 14$ days (the lookback):
\begin{equation}\label{eq:sigma}
  \sigma_{t,k} = \frac{1}{n}\sum_{i=1}^{n} m_{t-i,k} .
\end{equation}
Since the sum starts at $i=1$, $\sigma_{t,k}$ only uses days before $t$, so it is $\mathcal F_{t,0}$-measurable.
In the code this is a rolling mean over days followed by a shift of one day.

\begin{remark}[Why the band widens like $\sqrt{k}$]
$\sigma_{t,k}$ is a mean \emph{absolute} move, not a standard deviation. Suppose one-minute log returns
are i.i.d.\ $\mathcal N(0,s^2)$. Then $\ln(P_{t,k}/O_t) \sim \mathcal N(0, k s^2)$, and for small moves
\[
  \E[m_{t,k}] \approx \E\,|Z|\, s\sqrt{k} = \sqrt{\tfrac{2}{\pi}}\; s\sqrt{k}, \qquad Z\sim\mathcal N(0,1).
\]
So the threshold for a ``significant'' move grows like $\sqrt{k}$. This gives the concave, funnel-shaped
Noise Area in Figure~1 of the paper. A fixed threshold would be too loose in the morning and too tight
in the afternoon.
\end{remark}

\subsection{Boundaries with gap adjustment}
Let $C_{t-1}$ be the previous close and $D_t$ the cash dividend if $t$ is an ex-dividend date
($D_t = 0$ otherwise). The dividend-adjusted previous close is $\tilde C_{t-1} = C_{t-1} - D_t$.
With the volatility multiplier $\mathrm{VM}$ (paper: $\mathrm{VM}=1$) the boundaries are
\begin{align}
  \UB_{t,k} &= \max\!\big(O_t,\, \tilde C_{t-1}\big)\,\big(1 + \mathrm{VM}\,\sigma_{t,k}\big), \label{eq:ub}\\
  \LB_{t,k} &= \min\!\big(O_t,\, \tilde C_{t-1}\big)\,\big(1 - \mathrm{VM}\,\sigma_{t,k}\big). \label{eq:lb}
\end{align}
The Noise Area is the interval $\mathcal N_{t,k} = [\LB_{t,k}, \UB_{t,k}]$.

\begin{remark}[Gap adjustment]
After a gap down ($O_t < \tilde C_{t-1}$) the upper band is anchored at yesterday's close. A long signal
then requires the price to close the whole gap \emph{and} move an extra $\sigma_{t,k}$. The lower band
stays anchored at the open, so continuation of the gap is detected quickly. Gap ups work the same way
in the other direction. Without the dividend correction, the mechanical price drop on ex-dates would look
like a gap down.
\end{remark}

\subsection{Session VWAP}
With the typical price $\bar p_{t,j} = (H_{t,j}+L_{t,j}+C_{t,j})/3$,
\begin{equation}
  \VWAP_{t,k} = \frac{\sum_{j=0}^{k-1} \bar p_{t,j}\,V_{t,j}}{\sum_{j=0}^{k-1} V_{t,j}},
\end{equation}
which uses bars up to minute $k$ only, so it is $\mathcal F_{t,k}$-measurable.

%=======================================================================
\section{Signal and position (\texttt{backtest.py})}
%=======================================================================

\subsection{Decision times}
Positions can only change at the decision times
\[
  \mathcal K_t = \{\,30,\,60,\,\dots\,\} \cap \{1,\dots,N_t-1\},
\]
i.e.\ 10:00, 10:30, \dots, 15:30 on a full day. Let $x_{t,k} \in \{-1, 0, +1\}$ be the position
(short, flat, long) held after the decision at minute $k$. Each day starts flat, $x_{t,0} = 0$, and
between decisions the position is unchanged.

\subsection{Base model: stop at the opposite band}
At $k\in\mathcal K_t$, with $x^- = $ the current position:
\begin{equation}\label{eq:base}
  x_{t,k} =
  \begin{cases}
    +1 & \text{if } P_{t,k} > \UB_{t,k},\\
    -1 & \text{if } P_{t,k} < \LB_{t,k},\\
    x^- & \text{otherwise.}
  \end{cases}
\end{equation}
This rule has \emph{memory}: inside the Noise Area the previous position is kept. A long is only closed
(and immediately reversed) when the price crosses the \emph{lower} band.

\subsection{Extension 1: trailing stop at the current band and VWAP}
\begin{equation}\label{eq:vwap}
  x_{t,k} =
  \begin{cases}
    +1 & \text{if } P_{t,k} > \max(\UB_{t,k}, \VWAP_{t,k}),\\
    -1 & \text{if } P_{t,k} < \min(\LB_{t,k}, \VWAP_{t,k}),\\
    0  & \text{otherwise.}
  \end{cases}
\end{equation}
This rule is \emph{memoryless}. As soon as a long falls back below the upper band or below VWAP, it is
closed. This is exactly a trailing stop at $\max(\UB_{t,k},\VWAP_{t,k})$.

\subsection{Execution}
A decision at $(t,k)$ is executed at the open of the next bar,
\[
  \hat P_{t,k} = O_{t,k} \quad(\text{realised after } \mathcal F_{t,k}),
\]
and every open position is closed at the day's close $C_t$. In particular $x_{t,N_t} = 0$:
no overnight exposure.

%=======================================================================
\section{Position sizing}
%=======================================================================

The number of shares is fixed at the open of each day:
\begin{equation}\label{eq:shares}
  q_t = \left\lfloor \frac{A_{t-1}\, L_t}{O_t} \right\rfloor ,
\end{equation}
where $A_{t-1}$ is the equity at the end of day $t-1$ and $L_t$ is the leverage.

\paragraph{Full exposure (base, Extension 1).} $L_t = 1$.

\paragraph{Volatility targeting (Extension 2).} Let $r_s = C_s / C_{s-1} - 1$ be daily close-to-close returns and
\begin{equation}
  \hat\sigma_t = \sqrt{\frac{1}{n-1}\sum_{i=1}^{n}\big(r_{t-i} - \bar r_t\big)^2}, \qquad
  \bar r_t = \frac1n\sum_{i=1}^n r_{t-i},
\end{equation}
the realised daily volatility over the previous $n=14$ days ($\mathcal F_{t,0}$-measurable). Then
\begin{equation}\label{eq:lev}
  L_t = \min\!\left(L_{\max},\; \frac{\sigma^\star}{\hat\sigma_t}\right),
  \qquad \sigma^\star = 2\%,\quad L_{\max} = 4 .
\end{equation}
Exposure is reduced when the market is volatile and increased when it is calm. Since SPY's daily volatility
is typically around $1\%$, the typical leverage is about $2$.

%=======================================================================
\section{Profit and loss with transaction costs}
%=======================================================================

Let the day's round trips be $\ell = 1,\dots,R_t$, each with side $s_\ell\in\{\pm1\}$,
entry price $e_\ell$ and exit price $z_\ell$. Every entry, exit or reversal leg is one order of $q_t$
shares. Let $M_t$ be the number of orders. A reversal from $+1$ to $-1$ counts as two orders.

\paragraph{Gross and net P\&L.}
\begin{equation}
  \Pi^{\text{gross}}_t = \sum_{\ell=1}^{R_t} s_\ell\, q_t\,(z_\ell - e_\ell),
  \qquad
  \Pi_t = \Pi^{\text{gross}}_t - M_t\,\kappa(q_t),
\end{equation}
with the cost of one order
\begin{equation}
  \kappa(q) = \max\big(c_{\min},\; c\,q\big) + \delta\, q ,
\end{equation}
where $c = \$0.0035$ is the commission per share, $c_{\min} = \$0.35$ the minimum commission per order and
$\delta$ the slippage per share ($\$0.001$ in the paper scenario, $\$0.005$ in the conservative one).

\paragraph{Equity and returns.}
\begin{equation}
  A_t = A_{t-1} + \Pi_t, \qquad R^{\text{strat}}_t = \frac{\Pi_t}{A_{t-1}}, \qquad A_0 = \$100{,}000 .
\end{equation}

\begin{prop}[Break-even cost]
Let $\bar\pi$ be the average gross profit per share per round trip. Each round trip has two orders.
For large $q$ (where the minimum commission does not bind), the strategy is profitable on average iff
\[
  \bar\pi > 2\,(c + \delta) .
\]
\end{prop}
\begin{proof}
A round trip earns $q\,\bar\pi$ on average and costs $2\kappa(q) = 2q(c+\delta)$ when $cq\ge c_{\min}$.
\end{proof}
\noindent The paper reports $\bar\pi \approx \$0.09$, so the edge vanishes at roughly
$c+\delta \approx \$0.045$ per share and side. This is the reason the cost assumptions matter so much
for an intraday strategy.

%=======================================================================
\section{Performance metrics (\texttt{metrics.py})}
%=======================================================================

Let $R_1,\dots,R_T$ be daily returns, with $\bar R$ their mean and $s_R$ their sample standard deviation.
We use $252$ trading days per year.

\begin{align}
  \text{Total return} &= \prod_{t=1}^T (1+R_t) - 1, \\
  \text{Annualised return (CAGR)} &= \Big(\prod_{t=1}^T (1+R_t)\Big)^{252/T} - 1, \\
  \text{Annualised volatility} &= s_R \sqrt{252}, \\
  \text{Sharpe ratio} &= \frac{\bar R - r_f/252}{s_R}\,\sqrt{252} \quad (r_f = 0 \text{ in our tables}), \\
  \text{Max drawdown} &= \max_{t}\left(1 - \frac{E_t}{\max_{u\le t} E_u}\right), \qquad E_t = \prod_{u\le t}(1+R_u), \\
  \text{Hit ratio} &= \frac{\#\{t : R_t > 0\}}{\#\{t : R_t \neq 0\}} .
\end{align}

\begin{remark}[Why $\sqrt{252}$]
If daily returns are independent with variance $s^2$, the variance of a sum over 252 days is $252\,s^2$.
The standard deviation therefore scales with $\sqrt{252}$, while the mean scales with $252$. Hence the
Sharpe ratio scales with $252/\sqrt{252} = \sqrt{252}$.
\end{remark}

\begin{prop}[Leverage does not change the Sharpe ratio]
For $r_f = 0$ and any constant $\lambda>0$: $\mathrm{SR}(\lambda R) = \mathrm{SR}(R)$.
\end{prop}
\begin{proof}
$\overline{\lambda R} = \lambda \bar R$ and $s_{\lambda R} = \lambda s_R$, so the ratio is unchanged.
\end{proof}
\noindent Therefore strategies should be compared by their Sharpe ratio and not by total return.
Volatility targeting changes $L_t$ over time, so it \emph{can} change the Sharpe ratio. Most of the
increase in total return from Extension 2, however, comes from the average leverage.

\paragraph{Alpha and beta.} OLS regression of strategy returns on SPY returns,
\begin{equation}
  R^{\text{strat}}_t = \alpha + \beta\, R^{\text{SPY}}_t + \varepsilon_t,
  \qquad
  \hat{\bm\theta} = (X^\top X)^{-1} X^\top y,
\end{equation}
with $X = [\bm 1, R^{\text{SPY}}]$. Standard errors are
$\widehat{\mathrm{se}}(\hat\theta_i) = \sqrt{\hat s^2_\varepsilon\,[(X^\top X)^{-1}]_{ii}}$ with
$\hat s^2_\varepsilon = \hat\varepsilon^\top\hat\varepsilon/(T-2)$, and $t_i = \hat\theta_i/\widehat{\mathrm{se}}(\hat\theta_i)$.
We report the annualised alpha $252\,\hat\alpha$.

\paragraph{Statistical significance of the Sharpe ratio.}
The $t$-statistic of the mean daily return is $t = \bar R/(s_R/\sqrt T) = \mathrm{SR}_{\text{daily}}\sqrt{T}$.
With $T = 252\,Y$ days ($Y$ years) this gives
\begin{equation}
  t \approx \mathrm{SR}_{\text{ann}}\,\sqrt{Y}.
\end{equation}
For example, a Sharpe ratio of $1$ over $4$ test years gives $t\approx2$, which is only borderline significant.

%=======================================================================
\section{Evaluation protocol}
%=======================================================================

\begin{itemize}
  \item \textbf{Split.} The sample is divided chronologically into a training period $\mathcal T_{\text{train}}$
        (2016--2021) and a test period $\mathcal T_{\text{test}}$ (2022--today). Each period is simulated
        separately with $A_0 = \$100{,}000$.
  \item \textbf{No fitting on the test set.} All variants use the paper's parameters
        $(n, \mathrm{VM}, \sigma^\star, L_{\max}) = (14, 1, 2\%, 4)$. The grid
        $n\in\{7,14,30,60\}$, $\mathrm{VM}\in\{0.8,1.0,1.2,1.5\}$ is evaluated on $\mathcal T_{\text{train}}$ only,
        as a robustness check. A good result is a flat plateau of Sharpe ratios rather than one sharp peak.
  \item \textbf{No look-ahead.} $\sigma_{t,k}$, $\hat\sigma_t$, $\tilde C_{t-1}$ and $L_t$ are
        $\mathcal F_{t,0}$-measurable. $P_{t,k}$ and $\VWAP_{t,k}$ are $\mathcal F_{t,k}$-measurable.
        Execution happens at $O_{t,k}$, which lies after $\mathcal F_{t,k}$. These properties are checked by
        unit tests (\texttt{tests/test\_signals.py}).
  \item \textbf{Costs.} Two scenarios ($\delta = 0.001$ and $\delta = 0.005$), plus a sweep over
        $c+\delta \in [0.0035, 0.0235]$ in the test period.
\end{itemize}

%=======================================================================
\section{Summary of parameters}
%=======================================================================

\begin{center}
\begin{tabular}{lll}
\toprule
Symbol & Meaning & Value \\
\midrule
$n$ & lookback for $\sigma_{t,k}$ and $\hat\sigma_t$ & 14 days \\
$\mathrm{VM}$ & volatility multiplier & 1 \\
$\mathcal K_t$ & decision times & every 30 min from 10:00 \\
$\sigma^\star$ & daily volatility target & 2\% \\
$L_{\max}$ & leverage cap & 4 \\
$c$ & commission per share & \$0.0035 \\
$c_{\min}$ & minimum commission per order & \$0.35 \\
$\delta$ & slippage per share & \$0.001 / \$0.005 \\
$A_0$ & initial capital & \$100{,}000 \\
\bottomrule
\end{tabular}
\end{center}

\end{document}
